\pagebreak

\section{SPI Library File - spi.3pl}

    \index{library!spi.3pl}
    SPI is a simple, synchronous serial protocol used to control peripheral
    chips through a four-wire bus, at typical speeds of 1Mbit/s. SPI is more
    or less defined in Section 8 of the Motorola MC68HC11 Reference Manual.
    A subset of SPI is known as Microwire.

    This implementation provides separate modules for SPI master and slave. Each
    module can receive and transmit, although in many cases it will only do one or
    the other. The SPI master is the initiator of all transfers. It generates the
    synchronous clock ({\bf SCK}) and chip select ({\bf \_nSS}) signals, and drives
    the master-out-slave-in ({\bf MOSI}) data line. The slave drives the
    master-in-slave-out ({\bf MISO}) data line. 

    Often there will be one master and one slave. However, single-master,
    multi-slave situations are common. In this case, the master drives multiple
    chip-select signals ({\bf \_nSS}), one per slave, and {\bf SCK}, {\bf MOSI} and
    {\bf MISO} are common. Slaves must in these cases use tri-state or
    open-collector logic to drive MISO. Multi-master operation is possible, but
    requires separate arbitration arrangements between masters, usually performed
    at a software level.

    Peripheral chip manufacturers often alter aspects of the interface, for
    example, inverting signals, using wired-or or tri-state logic, omitting {\bf
    \_nSS}, using a common wire for {\bf MOSI}/{\bf MISO}, and so on. This
    implementation should still be useful in these cases, using some additional
    glue logic  to account for the differences.
    
    This library must be incorporated by using a {\bf module} statement -
    \begin{lstlisting}[gobble=4, language=threepl]
       module "spi.3pl"
    \end{lstlisting}
    


    \subsection{library modules}

        \subsubsection{spi\_master - SPI master transmitter/receiver}\label{sec:lmspimaster}

            {\bf Library: spi.3pl}

            {\bf Prototype:} \\
            \vsp
            \begin{lstlisting}[gobble=12, language=threepl]
               global module spi_master(
                   static(SCK, "log"),
                   static(MOSI, "log"),
                   static(_NSS, "log"),
                   queue(rx)
               <-
                   queue(tx),
                   value(MISO, "log", false),
                   value(timebase, "log", true),
                   log(chpa, false)
               ) {}
            \end{lstlisting}

            {\bf Output parameters:} \\
            {\bf SCK} is the output data clock.\\
            {\bf MOSI} is the master data output to a slave. \\
            {\bf \_NSS} is a slave chip select. \\
            {\bf rx} is the data output queue for the input data from MISO.

            {\bf Input parameters:} \\
            {\bf tx} is the data input queue for the output data to MOSI.\\
            {\bf MISO} is the master data input from a slave.\\
            {\bf timebase} is the .\\
            {\bf chpa} controls the clock sense.

            {\bf description:} \\
            SPI master transmitter/receiver. Data written to the tx queue is
            serialised and clocked out to the {\bf MOSI} (master-out slave-in)
            output in synchronism with the generated {\bf SCK} output.
            Simultaneously, input data on {\bf MISO} is sampled, deserialised and
            written to the {\bf rx} queue. {\bf \_NSS} goes low for the duration of
            the transmission. {\bf SCK}, {\bf MOSI}, {\bf MISO} and {\bf \_NSS} are
            normally routed to FPGA pins for connection to SPI peripherals.

            The data word size is the larger of that of the rx and tx queues.
            Generally these will be the same. Where the receive size is larger, 
            transmit words will be zero-padded at trailing end. Where the transmit
            size is larger, the only the trailing receive bits will be returned.

            Timing is regulated by the {\bf timebase} signal. Normally this is a stream
            of single-clock pulses at twice the SPI baud rate; SPI peripheral chips
            normally run at 1Mbaud or lower. If the timebase signal is tied true,
            the baud rate will be one quarter the prevailing clock rate, and if the
            timebase parameter is omitted, it will be one half. Note that SPI
            slaves may not work reliably at high speeds.

            Input {\bf chpa} is an immediate input which controls the phase of {\bf
            SCK}. When false (the default), data will be sampled on the rising edge
            of of {\bf SCK}. When true, the falling edge is used. (Note this is not
            the same as inverting {\bf SCK}).

            The master may be configured transmitter-only by omitting the {\bf rx}
            queue and {\bf MISO} input. This is a common situation. Similarly, the 
            master may be configured receiver-only by making the type of the {\bf
            tx} queue null. In this case (empty) writes to the {\bf tx} queue initiate
            receives.

            If, following a receive, the {\bf rx} queue is not available, the received
            data word is silently discarded.



        \subsubsection{spi\_slave - SPI slave transmitter/receiver}\label{sec:lmspislave}

            {\bf Library: spi.3pl}

            {\bf Prototype:} \\
            \vsp
            \begin{lstlisting}[gobble=12, language=threepl]
               global module spi_slave(
                   value(MISO, "log"),
                   queue(rx),
                   static(SOSI, "log")
               <-
                   value(tx),
                   value(SCK, "log", false),
                   value(MOSI, "log", false),
                   value(_NSS, "log", true),
                   log(chpa, false)
               ) {}
            \end{lstlisting}

            {\bf Output parameters:} \\
            {\bf MISO} is the data output to the master.\\
            {\bf rx} is the output queue for data received from the master.\\
            {\bf SOSI} is an optional output to a daisy-chained slave.

            {\bf Input parameters:} \\
            {\bf tx} is the data input for transmission to the master.\\
            {\bf SCK} is the input clock from the master.\\
            {\bf MOSI} is the data input from the master.\\
            {\bf \_NSS} is the slave chip select from the master.\\
            {\bf chpa} controls the clock sense.

            {\bf description:} \\
            SPI slave transmitter/receiver. When the connected SPI master initiates
            a transfer, the value of input {\bf tx} is sampled, serialised and
            clocked out to the {\bf MISO} (master-out slave-in) output in
            synchronism with the received SCK input. Simultaneously, input data on
            {\bf MOSI} is sampled, deserialised and written to the {\bf rx} queue.
            {\bf \_NSS} must be low throughout the operation.

            Input {\bf chpa} controls the phase of the expected {\bf SCK}. The meaning is
            as described in module spi\_master() ({\bf not implemented yet}).

            The data word length is determined by the connected SPI master, although
            generally the slave will match. However, the word length may be dynamic,
            especially in a multi-master system.

            The slave may be configured as a transmitter only by omitting the {\bf
            rx} queue and {\bf MOSI} input. Similarly, the slave may be configured
            as a receiver only by omitting the the {\bf tx} parameter.

            If, during a receive, the {\bf rx} queue is not available (i.e., full),
            the received data word is silently discarded.

            In the unlikely event slave daisy-chaining is desired, the {\bf SOSI}
            output can be connected to the {\bf MOSI} input of subsequent slaves.
            {\bf SOSI} is the  output of the recieve shift-register, and so is the
            same as the slave's {\bf MOSI} input but delayed by the bit-width of
            the receiver. {\bf This feature is untested.}




    \subsection{library procedures}
    
        None.
        


    \subsection{library functions}
    
        None.


    \subsection{library classes}
    
        None.
